%% Time-stamp: <2013-02-07 11:51:00 vk>
%%%% === Disclaimer: =======================================================
%% created by
%%      Karl Voit
%%
%% edited by 
%%  	Andreas Fleischhacker

\usepackage{siunitx}
\usepackage{amsmath,amssymb,amsfonts}
\graphicspath{{figures/}}

\newcommand{\myfig}[6]{
\begin{figure}[#6]
  \begin{center}
     \includegraphics[keepaspectratio,#2]{figures/#1}
     \caption[#4]{#3}
     \label{#5} %% NOTE: always label *after* caption!
  \end{center}
\end{figure}
}



\newcommand{\myfigpdftex}[6]{
\begin{figure}[#6]
  \def\svgwidth{#2\textwidth}
  \begin{center}
    \input{figures/#1.pdf_tex}
    \caption[#4]{#3}
    \label{#5} %% NOTE: always label *after* caption!
  \end{center}
\end{figure}
}
% example:
%\myfigpdftex{}%% filename in figures folder
%       {}%% maximum height, aspect ratio will be kept
%       {}%% caption
%       {}%% optional (short) caption for list of figures
%       {}%% label
%       {}%% figure position

\newcommand{\vect}[1]{\mathbf{#1}} % vector symbol: bold
\newcommand{\mat}[1]{\mathbf{#1}}  % matrix symbol: bold
\newcommand{\esym}[1]{\textit{#1}} % symbols without a kursiv font
\newcommand{\Epsilon}{\mathcal{E}} % big epsilon
\newcommand{\underconsideration}[2]{\left.#1\right\vert_{#2}}
\newcommand{\germanword}[1]{\foreignlanguage{ngerman}{#1}}
\newcommand{\germanwordQM}[1]{\glqq\foreignlanguage{ngerman}{#1}\grqq}
\newcommand{\myeqitem}[2]{#1 & -- &  #2 \\}
\newenvironment{myquote}{\begin{quote} \small}{\end{quote}}
%\newenvironment{myeqdes}
%	{\tabularx{\textwidth}{lcX}}
%	{\endtabularx}
	



% Using \myclone[42]{foobar} results the text foobar printed 42 times.
% But you can not only repeat text output with myclone. 
%
% Default argument
% for the optional parameter number of times (like 42 in the example above) 
% is set to two.
% 
%% \myclone[x]{text}
\newcounter{myclonecnt}
\newcommand{\myclone}[2][2]{%
  \setcounter{myclonecnt}{#1}%
  \whiledo{\value{myclonecnt}>0}{#2\addtocounter{myclonecnt}{-1}}%
}


